\scanpage{74}
\uline{Un théorème d'orthogonalité}

\deleted{$X$ schéma}
$A$ anneau noethérien $J$-adique (séparé et complet)
\[
X=\operatorname{Spec}A\qquad
Y=\operatorname{Spec}(A/J)=V(J)
\]
\deleted{$\unsure{G}$}
\[
U=X-Y
\]
Pour tout schéma fini étale $Z$ sur $Y$, $\widetilde Z$ est le schéma fini étale sur $X$ (unique à isom unique près) le prolongeant.
\margin{\unsure{Idem} pour faisceaux $\ell$-adiques \unsure{constants tordus}, $\ldots$}
$G,G'$ schémas \deleted{\unsure{lisses}} \inserted{en groupes comm. \deleted{\unsure{séparés}} de t.f. séparés} \inserted{plats} sur $X$ \deleted{séparés}
\[
\begin{array}{ll}
G_U=G|U&\text{est un schéma abélien }A\\
G'_U=G'|U&\text{\rule{5cm}{.3pt}}\quad A'
\end{array}
\]
\margin{$A$ et $A'$ \deleted{\unsure{duaux l'un de l'autre}} \inserted{accouplés}, i.e. on a un élém. de corr. alg. sur $A\times A'$.}
$\ell$ \deleted{\unsure{entier}} nb premier inv. sur $Y$, donc sur $X$.
\margin{\unsure{commencer par parler de} $\unsure{Z,Z'}$, \unsure{puis parler de} $\deleted{\unsure{G}}A,A'$}
$S\subset G_Y=G\times_XY$ est \inserted{plat sur $Y$ et} tel que ${}_{\ell^\nu}S$ \uline{fini} sur $Y$ pour tout $\nu>0$.
\deleted{$\unsure{S'\subset G'_Y=G'\times_XY}$ \unsure{avec} $\unsure{{}_{\ell^\nu}S'}$ \unsure{fini sur} $Y$}
\[
T'\subset G'_Y\qquad\text{un \uline{tore}}
\]
\deleted{\unsure{Alors}}
$\widetilde{T_\ell(S)}\deleted{\unsure{\subset}}$ est défini, et il est $\subset T_\ell(G)$
idem $\widetilde{T_\ell(T)}$ défini et $\subset T_\ell(G')$
\margin{\uline{NB} les ${}_{\ell^\nu}G$ sont étales de t.f. sép. sur $S$, et leur restriction à $Y$ contient les ${}_{\ell^\nu}S$, donc ${}_{\ell^\nu}G$ contient ${}_{\ell^\nu}\widetilde S$.}
\diagram{\begin{tikzcd}[row sep=small]
\widetilde{T_\ell(S)}|U\subset T_\ell(A)\\[-.4em]
\times\\[-.4em]
\widetilde{T_\ell(T)}|U\subset T_\ell(A')\arrow[d,"\varphi_\xi"',"\text{\unsure{grâce à}}\ \xi"]\\
T_\ell(\unsure{\mathbb G|\mathbf G}_{mU})=\widetilde{T_\ell(\unsure{\mathbb G|\mathbf G}_{mY})}|U
\end{tikzcd}}
\margin{induit $T_\ell(\unsure{S|\xi})$}
\uline{Théorème}
\[
\varphi_\xi\bigl(\widetilde{T_\ell(S)}|U,\widetilde{T_\ell(T)}|U\bigr)=0
\]
\uline{Dém} Il suffit de \uline{le} prouver en chaque pt $x\in U$.
Prenant un point \inserted{$y\in\overline{\{x\}}\cap Y$} \deleted{dans l'adhérence}, et dominant par un anneau de val. discret, on est ramené au cas où $X$ est un \uline{trait}, $Y$ son pt fermé, $U$ son pt gén. Il suffit d'ailleurs, dans cette réduction, de prendre pour $x$ un pt maximal de $U$. Supposons \deleted{que} d'abord qu'on
\mspage{73}
